%\pagebreak
\section{\kcode{safesync} Clause}
\label{sec:safesync}

\index{clauses!safesync@\kcode{safesync}}
\index{safesync clause@\kcode{safesync} clause}

Parallel regions in OpenMP are usually executed by a team of threads that
should be able to freely take \emph{divergent code paths} and synchronize with
each other using various point-to-point synchronization mechanisms. However, on
certain target architectures in which OpenMP threads may instead execute a
single stream of (possibly predicated) instructions in lock step, two threads
may not be able to synchronize with each other from logically divergent code
without a resulting deadlock. The \kcode{safesync(\plc{width})} clause may be
specified on a \kcode{parallel} construct to control which threads in a team
should be able to synchronize with each other from divergent code paths. The
effect is that the team is partitioned into ``progress groups'' of consecutive
threads (in terms of thread number) of size \plc{width}, except for the final
group which can have less than \plc{width} threads. Threads in the same
progress group may execute a single sequence of instructions, while threads in
different progress groups will execute distinct sequences of instructions. Thus,
two threads in different progress groups can safely synchronize from divergent
code paths.

If the \kcode{safesync} clause isn't present the behavior is as if the
\plc{width} is 1 for a \kcode{parallel} construct encountered on the host
device or when the \kcode{device_safesync} requirement has been specified for
that compilation unit. In other words, any two threads in the team can
synchronize from divergent code paths by default under those conditions.

The following example uses a ticket lock implementation in a \kcode{target}
region that requires thread synchronization from divergent code paths. Under
the progress guarantees for threads that are defined in the 6.0 specification,
this code can deadlock without an explicit use of the \kcode{safesync} clause
when executing on a non-host device. For example, on an NVIDIA GPU, the
implementation may map threads in the team to individual GPU threads in a warp
which does not support such synchronization between divergent threads. To make
this work, the \kcode{safesync} clause (with the omitted \plc{width}
defaulting to 1) can be specified. The example includes two variants of a
\kcode{parallel} directive, one with a \kcode{safesync} clause and one without,
via a metadirective that is conditioned on whether the device is a host device
or non-host device. Alternatively, the \kcode{device_safesync} requirement
could be specified with a \kcode{requires} directive, and this would require
\kcode{safesync} behavior as a default even on a non-host device.

\cexample[6.0]{safesync}{1}
\ffreeexample[6.0]{safesync}{1}

